\scanpage{73}\mspage{H}\deleted{\unsure{H}}
c'est un isomorphisme à cause du fait que
$\pi_1(X)\unsure{\to|-}\pi_1(Y)$ est un isomorphisme [grâce au
th. de pureté sur $Y$]
\margin{$Y$ \uline{régulier} \unsure{$[$|$\Gamma$}}
% slanted

Le deuxième homomorphisme s'écrit
\[
H^0(Y,\underline{O}_Y/\underline{O}_Y^{p^{\unsure{\nu|v}}})\to H^0(X,\underline{O}_X/\underline{O}_X^{p^\nu})
\]
[grâce au fait que $X$, $Y$ sont réduits],
ou aussi, supposant que $k$ est alg. clos
donc parfait, s'interprète comme
\margin{$Y$ \unsure{simple|intègre} sur $k$ (?)}
% slanted
\[
f^*:H^0(Y,\underline{O}_Y^{p^{-\nu}}/\underline{O}_Y)\to H^0(X,\underline{O}_X^{p^{-\nu}}/\underline{O}_X)
\]

Ce dernier est bijectif grâce à la
\uline{normalité} de $X$ et $Y$, car \deleted{\unsure{si $f$}} avec
les notations ci-dessus, on a un
diagramme commutatif
\diagram{\begin{tikzcd}[ampersand replacement=\&,column sep=large,row sep=small]
{}\&T(X)\arrow[dl,"g^*"]\\
T(U)\&T(Y)\arrow[l,"i^*","\sim"']\arrow[u,"f^*"']
\end{tikzcd}}
avec $g^*$ \uline{injective} [car $\underline{O}_X^{p^{-\nu}}/\underline{O}_X^{\unsure{?|2}}$ est \deleted{\unsure{l'}} 
\uline{sans torsion} ---], donc $f^*$ et $i^{*-1}g^*$ sont des isom
inverses l'un de l'autre.
\margin{$X$ \uline{\unsure{normale|normal ?|normal}}}
% slanted
% archive sheet number: 72
